\documentclass[12pt]{ximera}
\title{Exam 1 (2022)}
\begin{document}
\begin{abstract}
An archived exam on Chapter 1: full rankings by three methods, a Condorcet candidate, a truncated Borda count, missing Borda points, Arrow's Impossibility Theorem, and counting pairwise comparisons.
\end{abstract}
\maketitle

\section*{Exam 1 (2022)}

This exam allowed one handwritten sheet of notes.

\begin{problem}
Use the following preference schedule to find the full ranking of candidates A, B, C,
D, E by each method. (The Borda count is not required here.)
\begin{center}
\begin{tabular}{c|ccccc}
Number of votes & 12 & 10 & 9 & 6 & 2\\\hline
1st & E & C & B & A & E\\
2nd & A & B & D & D & C\\
3rd & C & A & C & B & A\\
4th & B & E & A & C & D\\
5th & D & D & E & E & B
\end{tabular}
\end{center}

\begin{prompt}
\textbf{(a)} Plurality: 1st $\answer{E}$, 2nd $\answer{C}$, 3rd $\answer{B}$,
4th $\answer{A}$, 5th $\answer{D}$

\begin{feedback}
First-place votes: E $14$, C $10$, B $9$, A $6$, D $0$.
\end{feedback}
\end{prompt}

\begin{prompt}
\textbf{(b)} Plurality-with-Elimination: 1st $\answer{B}$, 2nd $\answer{E}$,
3rd $\answer{C}$, 4th $\answer{A}$, 5th $\answer{D}$

\begin{feedback}
With $39$ voters a majority is $20$.
\begin{itemize}
\item Round 1: E $14$, C $10$, B $9$, A $6$, D $0$. Eliminate D.
\item Round 2: unchanged. Eliminate A; its $6$ voters rank B next (D is gone).
\item Round 3: B $15$, E $14$, C $10$. Eliminate C; its $10$ voters rank B next.
\item Round 4: B $25$, E $14$. B wins.
\end{itemize}
\end{feedback}
\end{prompt}

\begin{prompt}
\textbf{(c)} Pairwise Comparisons: 1st $\answer{C}$, 2nd $\answer{A}$, 3rd $\answer{B}$,
4th $\answer{E}$, 5th $\answer{D}$

\begin{feedback}
\[
\begin{array}{lll@{\qquad}lll}
\text{A vs B} & 20\text{--}19 & \text{A} & \text{B vs C} & 15\text{--}24 & \text{C}\\
\text{A vs C} & 18\text{--}21 & \text{C} & \text{B vs D} & 31\text{--}8 & \text{B}\\
\text{A vs D} & 30\text{--}9 & \text{A} & \text{B vs E} & 25\text{--}14 & \text{B}\\
\text{A vs E} & 25\text{--}14 & \text{A} & \text{C vs D} & 24\text{--}15 & \text{C}\\
\text{C vs E} & 25\text{--}14 & \text{C} & \text{D vs E} & 15\text{--}24 & \text{E}
\end{array}
\]
Points: C $4$, A $3$, B $2$, E $1$, D $0$.
\end{feedback}
\end{prompt}
\end{problem}

\begin{problem}
Is there a Condorcet candidate for the election in Problem 1?

\begin{multipleChoice}
\choice{Yes, A}
\choice{Yes, B}
\choice[correct]{Yes, C}
\choice{Yes, E}
\choice{No}
\end{multipleChoice}

\begin{feedback}
C wins all four of its comparisons. Notice that neither the Plurality winner (E) nor the
Plurality-with-Elimination winner (B) is the Condorcet candidate.
\end{feedback}
\end{problem}

\begin{problem}
A club elects its president using a truncated Borda count, as for the Heisman Trophy:
each voter lists a top three, worth $3$, $2$, and $1$ points. The candidates are Aaron,
Bella, Caroline, Damien, and Eva (A, B, C, D, E).
\begin{center}
\begin{tabular}{c|ccccc}
Number of votes & 6 & 5 & 3 & 3 & 2\\\hline
1st & A & D & C & E & D\\
2nd & C & E & A & C & B\\
3rd & B & C & B & A & A
\end{tabular}
\end{center}

\begin{prompt}
Find the points and the top three candidates.

A: $\answer{29}$ \quad B: $\answer{13}$ \quad C: $\answer{32}$ \quad D: $\answer{21}$
\quad E: $\answer{19}$

\smallskip
1st $\answer{C}$, 2nd $\answer{A}$, 3rd $\answer{D}$

\begin{feedback}
\[
\begin{aligned}
\text{A:}&\ 6\cdot3+3\cdot2+3\cdot1+2\cdot1=29\\
\text{B:}&\ 6\cdot1+3\cdot1+2\cdot2=13\\
\text{C:}&\ 6\cdot2+5\cdot1+3\cdot3+3\cdot2=32\\
\text{D:}&\ 5\cdot3+2\cdot3=21\\
\text{E:}&\ 5\cdot2+3\cdot3=19
\end{aligned}
\]
(Check: $19$ ballots at $6$ points each is $114$.) Caroline wins without a single
first-place vote from the largest group.
\end{feedback}
\end{prompt}
\end{problem}

\begin{problem}
An election with four candidates is decided by a modified Borda count with $7$ points
for 1st place, $3$ for 2nd, $2$ for 3rd, and $1$ for 4th. There are $20$ ballots. A gets
$93$ points, B $85$, and C $50$. How many points did D get? $\answer{32}$

\begin{feedback}
Each ballot awards $7+3+2+1=13$ points, for $20\cdot13=260$ in total. D has
$260-(93+85+50)=32$.
\end{feedback}
\end{problem}

\begin{problem}
Someone claims to have invented a new voting system, ``MyVote,'' for city council
elections. The details are a business secret, but they guarantee that MyVote satisfies
the Majority, Monotonicity, and Independence-of-Irrelevant-Alternatives criteria.

\begin{prompt}
\textbf{(a)} Which theorem shows that, even if these claims are true, MyVote must violate
another fairness criterion in some situations?

\begin{multipleChoice}
\choice[correct]{Arrow's Impossibility Theorem}
\choice{The Pythagorean Theorem}
\choice{Condorcet's Paradox}
\choice{May's Theorem}
\end{multipleChoice}
\end{prompt}

\begin{prompt}
\textbf{(b)} Which criterion must it violate?

\begin{multipleChoice}
\choice[correct]{The Condorcet criterion}
\choice{The Majority criterion}
\choice{The Monotonicity criterion}
\choice{None: MyVote could satisfy all four}
\end{multipleChoice}

\begin{feedback}
Arrow's Impossibility Theorem says no voting method can satisfy all four fairness
criteria in every election. So MyVote must sometimes violate the Condorcet criterion:
some elections will have a candidate who beats every other candidate head to head but
does not win under MyVote.
\end{feedback}
\end{prompt}
\end{problem}

\begin{problem}
Elections often have many obscure candidates. How many pairwise comparisons are needed
to evaluate an election with $28$ candidates? $\answer{378}$

\begin{feedback}
$\frac{28\cdot27}{2}=378$.
\end{feedback}
\end{problem}

\end{document}
